\documentclass[12pt,a4paper]{article}
\usepackage[utf8]{inputenc}
\usepackage[french]{babel}
\usepackage[T1]{fontenc}

\title{Dossier de Présentation : Intégration de l'entité Obscurus au sein de la M2L}
\author{Maël - BTS SIO SISR}
\date{2026}

\begin{document}

\maketitle

\section{Présentation du Contexte}
La Maison des Ligues de Lorraine (M2L) est un établissement qui fournit des espaces et des services aux structures sportives régionales[cite: 11, 12]. Dans le cadre de son développement, la M2L accueille une nouvelle structure partenaire : \textbf{Obscurus}.

\section{Problématiques Métiers}
L'intégration de cette nouvelle entité soulève deux défis majeurs pour le service informatique :

\subsection{Événementiel Hors-Les-Murs (Projet ASSIZ)}
L'entité Obscurus organise une manifestation sous chapiteau accueillant 2500 visiteurs à proximité de la M2L[cite: 33]. L'objectif est de déployer une infrastructure réseau temporaire capable de séparer les flux administratifs de l'organisation des flux publics destinés aux visiteurs[cite: 34, 41].

\subsection{Mobilité des Collaborateurs (Projet DYLAN)}
Les membres d'Obscurus utilisent régulièrement les espaces mutualisés de la M2L, notamment la salle de réunion commune[cite: 70]. Ils doivent pouvoir accéder à leurs ressources internes de manière sécurisée et automatique, quel que soit le port de connexion utilisé[cite: 71].

\section{Objectifs du Projet}
\begin{itemize}
    \item Assurer la confidentialité des données de l'organisation[cite: 49].
    \item Garantir un accès internet stable pour le public lors des événements[cite: 48].
    \item Automatiser la configuration réseau pour les postes nomades[cite: 71].
\end{itemize}

\end{document}